We report here the structure and intended contents of the empirical results for the graph-aware prompting pipeline, and provide qualitative discussion of the types of evidence and error modes we use to interpret outcomes. Because no experimental artifacts or retrieval excerpts were provided with the current authoring request, concrete numerical summaries and example excerpts are omitted from this text and will be supplied alongside the machine-readable result artifacts in the final submission\footnote{Numerical tables, reviewer scores, and example JSON outputs were not available to the authoring agent at the time of writing.}.

Quantitative summaries (planned presentation). Following the evaluation protocol described in Section 4, we organize quantitative reporting around three primary metrics: citation coverage, hallucination rate (unsupported assertions per section), and reviewer-style judgment scores. For each metric we will provide (a) per-domain aggregates, (b) per-section breakdowns (to expose within-document variance), and (c) confidence intervals or bootstrap estimates when sample sizes permit. All numerical tables will be accompanied by links or references to the corresponding JSON and LaTeX/bib artifacts so that individual section outputs and their provenance can be inspected programmatically.

Qualitative analysis and error-mode taxonomy. To complement aggregate statistics, we plan to present representative examples that illustrate successful evidence integration and common failure modes. The qualitative taxonomy used to organize examples will include at least the following categories: (1) evidence-aligned synthesis, where a generated claim is directly supported by injected snippets and an explicit citation; (2) citation omission, where a factual claim lacks any attached provenance; (3) citation mismatch, where an attached citation does not substantively support the assertion; and (4) structural inconsistency, where section-level claims conflict with material produced in dependent subsections. For each category we will include the JSON-constrained section output, the injected evidence snippets with provenance metadata, and a short annotator rationale.

Conservative interpretation and per-topic variability. Analyses will emphasize conservative interpretation: we report measured trends without extrapolating beyond the evaluated domains, and we highlight variability across topics and section types (e.g., methods vs. background). Where differences between the graph-aware and linear baselines are small or sensitive to retrieval quality, we will refrain from strong claims and instead describe the practical implications for editing and auditability.

Auditability and artifact release. To facilitate reproducibility and post-hoc auditing, all reported quantitative tables will reference the exact JSON section outputs and the assembled LaTeX + refs.bib used to compute metrics. These artifacts enable independent verification of citation coverage and unsupported-claim counts at the granularity of individual sentences and citations.

Summary. The final Results section will combine the quantitative tables described above with the qualitative taxonomy and representative exemplars, and will explicitly flag cases where observed differences depend on retrieval quality or dataset composition. Concrete values and exemplar snippets are intentionally omitted from this draft pending inclusion of the experimental artifacts and retrieval excerpts referenced in the evaluation protocol.